One has therefore proved that $R$ is a system of roots of $G'$ with respect to $T'$, and one finishes the proof without difficulty.
\begin{corollary}{4.3.2}\label{XXII.4.3.2}
Let $S$ be a scheme, $G$ a reductive $S$-group, $Q$ an invariant subgroup of multiplicative type of $G$. Then $Q$ is central in $G$, the quotient $G/Q$ is representable by a reductive $S$-group, and the canonical morphism $G\to G/Q$ is locally splittable for the \'etale topology (with radicial exponents equal to $1$).
\end{corollary}
The first assertion follows from Exp.~IX 5.5; the others are local for the \'etale topology and one is reduced to 4.3.1.
\begin{definition}{4.3.3}\label{XXII.4.3.3}
Let $G$ be a reductive $S$-group. One says that $G$ is \emph{adjoint} (resp. \emph{simply connected}) if for every $s\in S$, the type of $G$ at $s$ is given by an adjoint (resp. simply connected) root datum, i.e. (Exp.~XXI 6.2.6) such that $M$ is generated by $R$ (resp. $M^*$ generated by $R^*$).
\end{definition}
\begin{proposition}{4.3.4}\label{XXII.4.3.4}
\begin{enumeratei}
\item An adjoint (resp. simply connected) reductive group is semisimple.
\item If $T$ is a maximal torus of the adjoint (resp. simply connected) reductive group $G$ and if $\alpha$ is a root of $G$ with respect to $T$, then the infinitesimal root $\overline\alpha$ is nonzero on each fiber (resp. $\alpha^*$ is a monomorphism and the infinitesimal coroot $H_\alpha$ is nonzero on each fiber).
\end{enumeratei}
\end{proposition}
Indeed, (i) is trivial; (ii) is verified on the geometric fibers and follows at once from Exp.~XXI 6.2.8.
\begin{proposition}{4.3.5}\label{XXII.4.3.5}
\begin{enumeratei}
\item For the reductive group $G$ to be adjoint, it is necessary and sufficient that $\uCentr(G)=\{e\}_S$.
\item For every reductive group $G$, the quotient group $G/\uCentr(G)$ is an adjoint reductive group.
\end{enumeratei}
\end{proposition}
Indeed, one may assume $G$ split; then (i) is trivial (for $\uCentr(G)=\mathrm D_S(M/\Gamma_0(R))$), and (ii) follows from 4.3.1.
\begin{definition}{4.3.6}\label{XXII.4.3.6}
Let $G$ be a reductive $S$-group. One calls the \emph{adjoint group of $G$}, and denotes by $\operatorname{ad}(G)$, the group $G/\uCentr(G)$. One calls the \emph{radical of $G$}, and denotes by $\operatorname{rad}(G)$, the maximal torus (unique by Exp.~XII 1.12) of $\uCentr(G)$. One calls the \emph{semisimple group associated with $G$} the quotient $G/\operatorname{rad}(G)$.
\end{definition}
The preceding definitions are compatible with base change. If $s\in S$, $\operatorname{rad}(G)_{\overline s}$ is indeed the radical of $G_{\overline s}$ in the usual sense (Exp.~XIX 1.6).
\numpara{4.3.7}\label{XXII.4.3.7}
If $(G,T,M,R)$ is a split group, then $\operatorname{rad}(G)=\mathrm D_S(M/N)$, where $N=M\cap\mathcal V(R)$, hence the semisimple group associated with $G$ (as indeed the adjoint group of $G$) is endowed with a canonical splitting (4.3.1) and one has a diagram of split groups
\[
G\to\operatorname{ss}(G)\to\operatorname{ad}(G)
\]
corresponding to the canonical diagram of root data (Exp.~XXI 6.5.5)
\[
\operatorname{ad}(\mathscr R(G))\to\operatorname{ss}(\mathscr R(G))\to\mathscr R(G).
\]
\begin{remark}{4.3.8}\label{XXII.4.3.8}
Let $(G,T,M,R)$ be an adjoint (resp. simply connected) split $S$-group, $\Delta$ a system of simple roots of $R$. Then the family $\{\alpha\}_{\alpha\in\Delta}$, resp. $\{\alpha^*\}_{\alpha\in\Delta}$, induces an isomorphism
\[
T\isomto(\mathbf G_{m,S})^\Delta,\qquad\text{resp.}\quad(\mathbf G_{m,S})^\Delta\isomto T.
\]
Indeed, $M=\Gamma_0(R)$ (resp. $\ldots$) and $\Delta$ is a basis of the free abelian group $\Gamma_0(R)$ (Exp.~XXI 3.1.8).
\end{remark}
\begin{remark}{4.3.9}\label{XXII.4.3.9}
The radical of a reductive group is a \emph{characteristic} subgroup (i.e. stable under $\underline{\Aut}_{S\text{-gr.}}(G)$), in view of its definition.
\end{remark}

\sgasection{5}{Subgroups of type (R)}\label{XXII.sec.5}
We are particularly interested in reductive groups, but some of the results we are going to establish hold more generally for a larger class of groups: groups of type (RR).
\sgasection{5.1}{Groups of type (RR)}
\begin{definition}{5.1.1}\label{XXII.5.1.1}
Let $S$ be a scheme, $G$ an $S$-group scheme. One says that $G$ is \emph{of type (RR)} if it satisfies the following conditions:
\begin{enumeratei}
\item $G$ is smooth and of finite presentation over $S$ and has connected fibers.
\item $G$ possesses maximal tori locally for the (fpqc) topology.
\item For every $s\in S$, every maximal torus $T$ of $G_{\overline s}$ and every root of $G_{\overline s}$ with respect to $T_{\overline s}$ (Exp.~XIX 1.10), $\uLie(G_{\overline s})^\alpha$ has rank $1$ (as a vector space over $\kappa(\overline s)$).
\item For every $s\in S$ and every maximal torus $T$ of $G_{\overline s}$, denote by $R$ the set of roots of $G_{\overline s}$ with respect to $T$ and by $\Gamma_0(R)$ the subgroup of $\Hom_{\overline s\text{-gr.}}(T,\mathbf G_{m,\overline s})$ generated by $R$; then the content{\renewcommand{\thefootnote}{(22)}\nde{The content of the root $\alpha$ is the positive generator of the ideal $\{f(\alpha),\ f\in\Gamma_0(R)^*\}$ of $\ZZ$; it is the largest integer $c>0$ such that $\alpha/c\in\Gamma_0(R)$.}} of every root $\alpha\in R$ in the free abelian group $\Gamma_0(R)$ (which is thus an integer $>0$) is invertible on $S$.
\end{enumeratei}
\end{definition}
\begin{enoncerm}{Recall}{5.1.1.1}\label{XXII.5.1.1.1}
{\renewcommand{\thefootnote}{(23)}\nde{This reminder has been added.}}%
Recall that if $G$ is a smooth connected algebraic group over an algebraically closed field $k$, a \emph{Cartan subgroup} of $G$ is the centralizer of a maximal torus of $G$ (XII 1.0), and such a subgroup is \emph{smooth and connected}: for this, as well as for other characterizations of Cartan subgroups, see \textit{Bible}, \S~7.1, Th.~1 in the case $G$ affine and Exp.~XII Th.~6.6 in the general case. If $S$ is an arbitrary scheme and $G$ a smooth $S$-group of finite type, one calls a \emph{Cartan subgroup} of $G$ a smooth $S$-subgroup $C$ of $G$ such that, for every $s\in S$, $C_{\overline s}$ is a Cartan subgroup of $G_{\overline s}$ (XII Def.~3.1).
\end{enoncerm}
\begin{remark}{5.1.2}\label{XXII.5.1.2}
a) By Exp.~XII 7.1 (where the hypothesis $G$ separated holds here since $G$ has connected fibers, \cf Exp.~VIB 5.5), (i) and (ii) imply that $G$ possesses maximal tori (resp. Cartan subgroups) locally for the \'etale topology, conjugate locally for the \'etale topology.

b) The Cartan subgroups of $G$ have connected fibers (Exp.~XII 6.6).

c) If $G$ is affine over $S$, (i) and (ii) are equivalent respectively to
\begin{quote}
(i$'$) $G$ is smooth over $S$ and has connected fibers.

(ii$'$) The reductive rank of the fibers of $G$ is locally constant (Exp.~XII 1.7).
\end{quote}
d) By Exp.~XII 8.8 (c) and (d), $G$ possesses a reductive center $Z$ and for every $s\in S$ one has, with the notations of (iv),{\renewcommand{\thefootnote}{(24)}\nde{The original indicated ``under the conditions of (iv)'', but the last condition of (iv) does not seem to be used here.}} $Z_{\overline s}=\bigcap_{\alpha\in R}\Ker(\alpha)$, whence
\[
\Hom_{\overline s\text{-gr.}}((T/Z)_{\overline s},\mathbf G_{m,\overline s})\simeq\Gamma_0(R).
\]
e) Condition (iv) holds in particular in the following two cases:
\begin{quote}
(1) $S$ has characteristic $0$;

(2) every root is an indivisible element of the group generated by the roots.
\end{quote}
f) Being of type (RR) is stable under base change and is local for the (fpqc) topology.
\end{remark}
From remarks (c) and (e), one obtains at once the
\begin{proposition}{5.1.3}\label{XXII.5.1.3}
A reductive group is of type (RR).
\end{proposition}
\begin{proposition}{5.1.4}\label{XXII.5.1.4}
Let $S$ be a scheme, $G$ an $S$-group of type (RR), $Q$ a central subgroup of $G$ of finite presentation over $S$ such that the quotient $G/Q$ is representable (for example $G$ affine over $S$ and $Q$ of multiplicative type (IX 2.3), or $S$ artinian (VIA 3.3.2)); then $G/Q$ is an $S$-group of type (RR).
\end{proposition}
Indeed $G/Q$ is smooth over $S$ (Exp.~VIB 9.2), of finite presentation over $S$ (Exp.~V 9.1) and has connected fibers, so condition (i) holds. On the other hand, condition (ii) follows from Exp.~XII 7.6; it remains to verify conditions (iii) and (iv).

Write $G'=G/Q$, let $u:G\to G'$ be the canonical morphism, $T'=u(T)$ the maximal torus of $G'$ which is the image of $T$ (\cf Exp.~XII 7.1 (e)); for each $\alpha\in R$, denote again by $\alpha$ the character of $T'$ defined by $\alpha$ (one has $Q\cap T\subset\bigcap_{\alpha\in R}\Ker(\alpha)$ by 5.1.2 (d)). Let us first prove:
\begin{lemma}{5.1.5}\label{XXII.5.1.5}
Under the conditions of 5.1.4, let $T=\mathrm D_S(M)$ be a trivialized maximal torus of $G$, assume that the decomposition of $\mathfrak g=\uLie(G)$ under $\Ad(T)$ is of the form
\[
\mathfrak g=\mathfrak g^0\oplus\coprod_{\alpha\in R}\mathfrak g^\alpha,\qquad R\subset M-\{0\},
\]
where for every $s\in S$, $\mathfrak g^\alpha(s)\ne0$ for $\alpha\in R$ (so $\mathfrak g^\alpha$ is an invertible $\OS$-module for every $\alpha\in R$ and $R$ is the set of roots of $G_{\overline s}$ with respect to $T_{\overline s}$ for every $s\in S$).

Then the Lie algebra $\mathfrak g'$ of $G'$ decomposes under $\Ad(T')$ as follows:
\[
\mathfrak g'=\mathfrak g'^0\oplus\coprod_{\alpha\in R}\mathfrak g'^\alpha
\]
and $\uLie(u)$ induces an isomorphism of $\mathfrak g^\alpha$ onto $\mathfrak g'^\alpha$.
\end{lemma}
Indeed, let $p=\uLie(u):\mathfrak g\to\mathfrak g'$. One has at once $p(\mathfrak g^\alpha)\subset\mathfrak g'^\alpha$ for every $\alpha\in R$, and $p(\mathfrak g^0)\subset\mathfrak g'^0$. Since
\[
\Ker(p)=\uLie(Q)\subset\uLie(\uCentr_G(T))=\mathfrak g^0,
\]
$p$ induces a monomorphism of $\mathfrak g^\alpha$ into $\mathfrak g'^\alpha$, for every $\alpha\in R$.

To prove the lemma, it suffices to do so when $S$ is the spectrum of an algebraically closed field, and by the preceding remarks it then suffices to prove that $\operatorname{rg}(\mathfrak g')=\operatorname{rg}(\mathfrak g'^0)+\operatorname{Card}(R)$. Now put $C=\uCentr_G(T)$, $C'=\uCentr_{G'}(T')$; by Exp.~XII 7.1 (e), $u$ induces a faithfully flat morphism $C\to C'$ with kernel $Q$. One therefore has
\[
\dim C'+\dim Q=\dim C.
\]
But $G$, $G'$, $C$ and $C'$ are smooth, hence
\[
\begin{aligned}
\dim G=\operatorname{rg}(\mathfrak g)&=\operatorname{rg}(\mathfrak g^0)+\operatorname{Card}(R)=\dim C+\operatorname{Card}(R)\\
&=\dim Q+\dim C'+\operatorname{Card}(R)\\
&=\dim Q+\operatorname{rg}(\mathfrak g'^0)+\operatorname{Card}(R),\\
\operatorname{rg}(\mathfrak g')&=\dim G'=\dim G-\dim Q,
\end{aligned}
\]
which implies
\[
\operatorname{rg}(\mathfrak g')=\operatorname{rg}(\mathfrak g'^0)+\operatorname{Card}(R),
\]
that is, the desired relation.

Let us return to the proof of 5.1.4; one may assume that $S$ is the spectrum of an algebraically closed field. Let $T$ be a maximal torus of $G$; applying 5.1.5, one has at once (iii) and (iv) for $G/Q$.

To use the preceding proposition, let us introduce a definition:
\begin{definition}{5.1.6}\label{XXII.5.1.6}
One says that the $S$-group scheme $G$ is \emph{of type (RA)} if it is of type (RR) and satisfies in addition condition (iv$'$) (stronger than (iv)):

(iv$'$) For every $s\in S$ and every maximal torus $T$ of $G_{\overline s}$, every root of $G_{\overline s}$ with respect to $T$ has a content in $\Hom_{\overline s\text{-gr.}}(T,\mathbf G_{m,\overline s})$ which is invertible on $S$.
\end{definition}
\begin{remarks}{5.1.7}\label{XXII.5.1.7}
a) An adjoint reductive $S$-group is of type (RA).

b) If $S$ has characteristic $0$, every group of type (RR) is of type (RA).

c) Being of type (RA) is stable under base change and is local for the (fpqc) topology.
\end{remarks}
The preceding remark (a) generalizes as:
\begin{proposition}{5.1.8}\label{XXII.5.1.8}
Let $S$ be a scheme, $G$ an $S$-group of type (RR), $Z$ its reductive center; assume the quotient $G/Z$ representable (for example $G$ affine over $S$, or $S$ artinian); then $G/Z$ is of type (RA).
\end{proposition}
Indeed, this follows at once from 5.1.4, 5.1.5 and 5.1.2 (d).
\begin{remark}{5.1.9}\label{XXII.5.1.9}
If $G$ is of type (RR) and if $T$ is a maximal torus of $G$, one may apply Exp.~XIX 6.3. In particular $W_G(T)$ is \'etale, quasi-finite and separated over $S$.
\end{remark}

\sgasection{5.2}{Subgroups of type (R)}
\begin{definition}{5.2.1}\label{XXII.5.2.1}
Let $S$ be a scheme, $G$ a smooth $S$-group scheme of finite presentation with connected fibers,{\renewcommand{\thefootnote}{(25)}\nde{The hypothesis that $G$ (resp. $H$) is of finite presentation over $S$ holds automatically since $G$ (resp. $H$), being smooth over $S$ and with connected fibers, is quasi-compact and separated over $S$ (VIB 5.5), hence of finite presentation over $S$.}} $H$ a group subscheme of $G$. One says that $H$ is \emph{of type (R)} if:
\begin{enumeratei}
\item $H$ is smooth, of finite presentation over $S$ and has connected fibers.{\renewcommand{\thefootnote}{(25)}\footnotemark[26]}
\item For every $s\in S$, $H_{\overline s}$ contains a Cartan subgroup of $G_{\overline s}$.
\end{enumeratei}
This notion is stable under base change and local for the (fpqc) topology.
\end{definition}
\begin{enoncerm}{Recall}{5.2.2}\label{XXII.5.2.2}
(\cf Exp.~XII 7.9) Under the preceding conditions:

a) $H=\uNorm_G(H)^0$.

b) If $G$ contains Cartan subgroups (resp. maximal tori) locally for the \'etale topology, the same holds for $H$, and the Cartan subgroups (resp. maximal tori) of $H$ are Cartan subgroups (resp. maximal tori) of $G$.
\end{enoncerm}
\begin{examples}{5.2.3}\label{XXII.5.2.3}
a) \emph{Borel subgroups}: a Borel subgroup of $G$ is a subgroup $H$ of type (R) whose geometric fibers are Borel subgroups of those of $G$.{\renewcommand{\thefootnote}{(26)}\nde{It amounts to the same to say that $H$ is a smooth subgroup of $G$, each of whose geometric fibers $H_{\overline s}$ is a Borel subgroup of $G_{\overline s}$ (since every Borel subgroup of $G_{\overline s}$ is connected and contains a Cartan subgroup of $G_{\overline s}$).}}

b) \emph{Parabolic subgroups}: a parabolic subgroup of $G$ is a subgroup of type (R) whose geometric fibers contain Borel subgroups.
\end{examples}
Other examples are given by the following propositions.
\begin{proposition}{5.2.4}\label{XXII.5.2.4}
Let $G$ be as in 5.2.1, $K\subset H$ two group subschemes of $G$, $H$ being assumed of type (R). Then $K$ is a subgroup of type (R) of $H$ if and only if it is a subgroup of type (R) of $G$.
\end{proposition}
{\renewcommand{\thefootnote}{(27)}\nde{What follows has been expanded.}}%
Indeed, let $s\in S$. Since $H$ is of type (R), every maximal torus of $H_{\overline s}$ is a maximal torus of $G_{\overline s}$, and therefore the same holds for Cartan subgroups.
\begin{proposition}{5.2.5}\label{XXII.5.2.5}
Let $G$ be as in 5.2.1, $T$ a maximal torus of $G$, $Q$ a subtorus of $T$, $Z=\uCentr_G(Q)$. If $H$ is a subgroup of type (R) of $G$ containing $T$, then $H\cap Z$ is a subgroup of type (R) of $Z$.
\end{proposition}
Recall first that $Z$ is a closed group subscheme of $G$, of finite presentation (Exp.~XI 6.11), with connected fibers (Exp.~XII 6.6), smooth over $S$ (Exp.~XI 2.4), hence satisfies the conditions imposed in definition 5.2.1. Likewise, $H\cap Z$ is of finite presentation, smooth and with connected fibers (for $H\cap Z=\uCentr_H(Q)$); moreover $H\cap Z\supset\uCentr_G(T)$.
\begin{proposition}{5.2.6}\label{XXII.5.2.6}
Let $S$ be a scheme, $G$ an $S$-group of type (RR) (resp. (RA)), $H$ a subgroup of type (R) of $G$. Then $H$ is an $S$-group of type (RR) (resp. (RA)).
\end{proposition}
Indeed (i) is clear, (ii) follows from 5.2.2 (b), (iii) and (iv) (resp. (iv$'$)) are to be verified when $S$ is the spectrum of an algebraically closed field. Then $H$ contains a maximal torus $T$ of $G$ (and hence also $C=\uCentr_G(T)${\renewcommand{\thefootnote}{(28)}\nde{By hypothesis, $H$ contains $\uCentr_G(T')$ for some maximal torus $T'$ of $G$; then $T$ and $T'$ are conjugate in $H$, hence $H$ also contains $C=\uCentr_G(T)$.}}) and the assertions to be proved follow at once from:
\begin{lemma}{5.2.7}\label{XXII.5.2.7}
Let $S$ be a scheme, $G$ an $S$-group of type (RR), $T$ a maximal torus of $G$ endowed with a trivialization $T\simeq\mathrm D_S(M)$; assume that
\[
\mathfrak g=\mathfrak g^0\oplus\coprod_{\alpha\in R}\mathfrak g^\alpha
\]
(the $\mathfrak g^\alpha$ being then invertible $\OS$-modules).

Let $H$ be a subgroup of type (R) containing $C=\uCentr_G(T)$ (i.e. containing $T$). Then $\mathfrak h=\uLie(H/S)$ is locally over $S$ of the form
\[
\mathfrak g^0+\coprod_{\alpha\in R'}\mathfrak g^\alpha=\mathfrak g_{R'};
\]
more precisely, let for each $s\in S$, $R'(s)=\{\alpha\in R\mid\mathfrak g^\alpha(s)\subset\mathfrak h(s)\}$. Then $R'(s)$ is a locally constant function of $s$; if $U$ is an open subset of $S$ on which $R'(s)=R'$, one has
\[
\mathfrak h_U=\mathfrak g_U^0\oplus\coprod_{\alpha\in R'}\mathfrak g_U^\alpha.
\]
\end{lemma}
Indeed, $\mathfrak h$ is a submodule of $\mathfrak g$, locally a direct summand, containing $\mathfrak g^0$ and stable under $T$.

\sgasection{5.3}{Strict transporter of two subgroups of type (R). Applications}
\begin{enoncerm}{Recall}{5.3.0}\label{XXII.5.3.0}
{\renewcommand{\thefootnote}{(29)}\nde{This reminder, which is used in the proof of 5.3.1 and 5.3.4, has been added.}}%
Let $S$ be a scheme, $G$ a smooth $S$-group, $\mathfrak g=\uLie(G/S)$ and $\mathfrak h$ an $\OS$-submodule of $\mathfrak g$ locally a direct summand. The $\OS$-algebra $\mathcal A=\operatorname{Sym}(\omega^1_{G/S})$ is locally free, hence the $S$-scheme $\uLie(G/S)=\mathrm W(\mathfrak g)=\Spec(\mathcal A)$ is essentially free in the sense of Exp.~VIII, 6.1. Since $\mathrm W(\mathfrak h)$ is a closed subscheme of $\uLie(G/S)$, of finite presentation over $\uLie(G/S)$, $N=\uNorm_G(\mathfrak h)$ is representable by a closed group subscheme of $G$, of finite presentation over $G$, by Exp.~VIII, 6.5 (a). (See also additions 6.2.3 and 6.2.4 (a) in Exp.~VIB.) On the other hand, by Exp.~II 5.3.1, one has $\uLie(N/S)=\uNorm_{\Lie(G/S)}(\mathfrak h)$.

Finally, by Exp.~VIB 3.10, if $N$ is smooth over $S$ at the points of the identity section, then the subgroup functor $N^0$ (defined in VIB 3.1) is represented by an open group subscheme of $N$, smooth over $S$.
\end{enoncerm}
